\documentclass[10pt,a4paper]{article}

% Codificacao e fontes vetoriais (sem serrilhado)
\usepackage[utf8]{inputenc}
\usepackage[T1]{fontenc}
\usepackage{lmodern}

% Layout da pagina
\usepackage[margin=1.8cm,top=1.4cm,bottom=1.4cm]{geometry}
\usepackage{setspace}
\setstretch{1.0}
\setlength{\parindent}{1cm}
\setlength{\parskip}{0pt}

% Hiperlinks
\usepackage{hyperref}
\hypersetup{colorlinks=true,linkcolor=black,filecolor=black,urlcolor=blue,citecolor=black}
\usepackage{url}

% Cabecalho e rodape
\usepackage{fancyhdr}
\pagestyle{fancy}
\fancyhf{}
\fancyhead[C]{\footnotesize Universidade Federal de Sao Carlos -- Departamento de Estatistica -- Mineracao de Dados}
\fancyfoot[C]{\thepage}
\renewcommand{\headrulewidth}{0.4pt}

\begin{document}

\thispagestyle{fancy}

\begin{center}
{\Large \textbf{Universidade Federal de Sao Carlos (UFSCar)}}\\[0.1cm]
{\large \textbf{Departamento de Estatistica}}\\[0.15cm]
{\large \textsc{Mineracao de Dados}}\\[0.3cm]
\rule{\textwidth}{0.8pt}\\[0.15cm]
{\LARGE \textbf{Projeto de Interpretabilidade de Modelos}}\\[0.1cm]
{\large \textbf{Projeto Mineracao -- Grupo 5}}\\[0.2cm]
\rule{\textwidth}{0.8pt}
\end{center}

\vspace{0.1cm}

\noindent\textbf{Integrantes:}
\begin{center}
\begin{tabular}{|l|l|}
\hline
\textbf{Nome} & \textbf{RA} \\
\hline
Arthur Pereira Gon & 811464 \\
Giovanni Roberti & 811676 \\
Igor Figueiredo & 791420 \\
Lucas Yukio & 813650 \\
Murilo Cassiavilani & 812067 \\
\hline
\end{tabular}
\end{center}

\section*{Tema e Modalidade}

\textbf{Tema sorteado:} 5 -- Interpretabilidade de Modelos. \textbf{Modalidade:} A -- Projeto aplicado de mineracao de dados.

\section*{Resumo}

Este projeto investiga a interpretabilidade de modelos preditivos aplicada aos cinco paradigmas de modelagem da disciplina. Para cada paradigma, foram selecionados datasets reais, ajustados modelos representativos e aplicadas tecnicas especificas de interpretacao para responder: por que o modelo prediz o que prediz?

\section*{Metodologia}

O desenvolvimento foi realizado em Python com scikit-learn para modelagem, utilizando validacao cruzada e divisoes temporais quando apropriado. As tecnicas de interpretabilidade empregadas incluem: coeficientes regularizados (Lasso), importancia por permutacao, Partial Dependence Plots (PDP) e SHAP (SHapley Additive exPlanations).

\section*{Paradigmas Implementados}

\textbf{1. Regressao -- Predicao de Salarios NBA.} Dataset: NBA Player Salaries 2022-23 (467 jogadores). Modelos: OLS, Ridge, Lasso, Random Forest. Variavel resposta: salario anual (log-transformado). Melhor resultado: Random Forest (R2 = 0.71). Interpretabilidade: coeficientes padronizados revelaram que minutos jogados e Box Plus/Minus sao os preditores mais importantes. Fonte: \url{https://www.kaggle.com/datasets/jamiewelsh2/nba-player-salaries-2022-23-season}.

\textbf{2. Classificacao -- Vitoria do Mandante na Premier League.} Dataset: Premier League Stats 2006-2018 (4.560 partidas). Modelos: Regressao Logistica com Lasso, Random Forest. Variavel resposta: vitoria do mandante (binaria). Divisao temporal: treino 2006-2015, teste 2016-2018. Melhor resultado: Logistica Lasso (AUC-ROC = 0.77). Interpretabilidade: o Lasso selecionou apenas 3 variaveis (diff\_wins, diff\_losses, diff\_att\_ibox\_goal), revelando que a forma do time e o unico preditor realmente importante. Fonte: \url{https://www.kaggle.com/datasets/zaeemnalla/premier-league}.

\textbf{3. Agrupamento -- Perfis de Jogadores de Tenis.} Dataset: Tennis Player Data (462 jogadores ATP). Algoritmos: K-Means (K = 2, escolhido via Silhouette = 0.24) e Agrupamento Hierarquico (Ward). Features: 12 metricas de saque e devolucao. Clusters identificados: (0) Saqueadores/Agressivos (42 jogadores, ex: Dimitrov, Hurkacz) e (1) Devolvedores/Baseline (56 jogadores, ex: Carreno Busta, De Minaur). Interpretabilidade: analise dos centroides, graficos radar e jogadores representativos validam a dicotomia classica do tenis. Fonte: \url{https://www.kaggle.com/datasets/mohammadkumail110/tennis-player-data-serve-return-and-raw}.

\section*{Paradigmas Pendentes}

\textbf{4. Mineracao de Textos:} Modelagem de topicos (LDA) e classificacao de documentos, com interpretacao por inspecao de palavras relevantes.

\textbf{5. Deteccao de Anomalias:} Isolation Forest e LOF, com interpretacao das variaveis que mais se desviam do perfil medio.

\section*{Implementacao}

Desenvolvimento em Python com scikit-learn, pandas, matplotlib e seaborn. Todo o codigo esta organizado em pipelines reprodutiveis, documentando desde a importacao dos dados brutos ate a geracao das figuras e tabelas finais. Repositorio: \url{https://github.com/ArthurGoon/Projeto-Mineracao-de-Dados}.

\section*{Referencias}

\begin{enumerate}
    \setlength{\itemsep}{0pt}
    \setlength{\parsep}{0pt}
    \item Welsh, J. NBA Player Salaries (2022-23 Season). Kaggle, 2023.
    \item Kumail, M. Tennis Player Data (Serve, Return and Raw). Kaggle, 2024.
    \item Zaeem, N. Premier League Stats. Kaggle, 2019.
    \item Tibshirani, R. Regression Shrinkage and Selection via the Lasso. JRSS, 1996.
    \item Breiman, L. Random Forests. Machine Learning, 2001.
    \item Lundberg, S. M.; Lee, S. I. A Unified Approach to Interpreting Model Predictions. NeurIPS, 2017.
\end{enumerate}

\end{document}
